\section{Methods}
\label{sec:methods}

\subsection{System architecture}
\label{sec:architecture}

\bsi\ is organised as a layered Python package with dependencies enforced
to point in a single direction, shown in Figure~\ref{fig:architecture}. A
\texttt{core} layer implements every nucleotide and protein calculation
against the Python standard library only, with no file, network, or
graphical-interface access. A \texttt{models} layer defines typed,
JSON-serialisable result dataclasses shared by every layer above it,
replacing the pattern -- present in the initial prototype -- of returning
pre-formatted display strings directly from analysis functions, which
precludes batch processing, machine-readable export, and independent testing
of computed values versus displayed wording. A \texttt{services} layer
provides HTTP transport (with pluggable backends, Section~\ref{sec:testing}),
a retry policy, an on-disk response cache, and structured logging. A
\texttt{structures} layer implements one client class per external resource
(RCSB, AlphaFold DB, ESM Atlas, UniProt) plus the validation logic of
Section~\ref{sec:validation} and a manager that sequences providers
according to configurable fallback order. A \texttt{workflows} layer
implements batch analysis and reproducible project management
(Section~\ref{sec:reproducibility}). An \texttt{io} layer implements CSV,
TSV, JSON, and self-contained HTML export. Finally, a command-line interface
and a Tkinter graphical interface both consume the same underlying API,
described in Section~\ref{sec:implementation}.

\begin{figure}[htbp]
\centering
\begin{tikzpicture}[
    layer/.style={rectangle, rounded corners, draw=black!60, minimum width=9.6cm,
                  minimum height=1.05cm, align=center, font=\small},
    arrow/.style={-{Stealth[length=2.2mm]}, thick, gray!70}
  ]
  \node[layer, fill=blue!8]  (gui)  {\textbf{Presentation} --- \texttt{gui/} (Tkinter), \texttt{cli.py}};
  \node[layer, fill=blue!14, below=0.28cm of gui]  (wf)   {\textbf{Orchestration} --- \texttt{workflows/} (batch runner, projects)};
  \node[layer, fill=blue!20, below=0.28cm of wf]   (str)  {\textbf{Domain services} --- \texttt{structures/} (providers, validation, manager)};
  \node[layer, fill=blue!26, below=0.28cm of str]  (svc)  {\textbf{Infrastructure} --- \texttt{services/} (transport, retry, cache, logging)};
  \node[layer, fill=blue!32, below=0.28cm of svc]  (core) {\textbf{Pure computation} --- \texttt{core/} and \texttt{models/} (stdlib only)};

  \draw[arrow] (gui.south) -- (wf.north);
  \draw[arrow] (wf.south) -- (str.north);
  \draw[arrow] (str.south) -- (svc.north);
  \draw[arrow] (svc.south) -- (core.north);

  \node[right=0.15cm of gui, font=\scriptsize, text width=2.6cm, align=left, text=gray!60]
       {depends on layer below};
\end{tikzpicture}
\caption{\bsi\ layer architecture. Dependencies point downward only;
\texttt{core} imports nothing outside the Python standard library, which is
verified by a dedicated continuous-integration job that installs the package
with no optional extras and fails if a third-party import is reachable from
\texttt{core}.}
\label{fig:architecture}
\end{figure}

Three design decisions follow directly from the motivating problem in
Section~\ref{sec:introduction}. First, results are dataclasses rather than
formatted strings, so that the value a validation check acts upon and the
text a user reads are produced from the same computation and can be tested
independently; the wording shown in the interfaces, including every warning
that accompanies a non-exact structure match, is itself covered by tests
(Section~\ref{sec:testing}). Second, sequence-to-structure validation
(Section~\ref{sec:validation}) is implemented as a single function called
uniformly by the manager after every provider returns a result, rather than
being available as an opt-in utility a caller might omit -- the failure mode
in Section~\ref{sec:introduction} arose precisely because validation was
optional in the initial prototype. Third, the HTTP transport is an
interface with three implementations (a \texttt{requests}-based transport,
a standard-library \texttt{urllib} transport used when \texttt{requests} is
not installed, and a scripted fake transport used only in tests), which is
what makes deterministic fault-injection testing possible
(Section~\ref{sec:fault-injection}): a test can assert that two HTTP 504
responses followed by a 200 response produce exactly one recovered result,
on every machine, without depending on a real service failing at a
convenient moment.

\subsection{DNA and protein sequence analysis}
\label{sec:core-calculations}

\paragraph{Nucleotide analysis.}
Input sequences are accepted in FASTA or raw form; header lines, whitespace,
and non-nucleotide punctuation copied incidentally from database records
(digits, alignment gaps) are ignored, while genuine unexpected symbols are
reported rather than silently discarded. Reported nucleotide measures
include base composition and GC content (IUPAC ambiguity codes excluded from
both numerator and denominator, so that a sequence with unresolved bases is
not artificially reported as GC-poor), reverse complementation, transcription,
six-frame translation, open-reading-frame detection with a configurable
minimum length and choice of NCBI genetic code table (standard, vertebrate
mitochondrial, mould/protozoan mitochondrial, and bacterial/archaeal tables
are bundled), and IUPAC degenerate-motif search on both strands.
Ambiguity codes are retained in place during translation rather than
stripped before codon splitting: a codon containing an unresolved base
translates to an unspecified residue (\texttt{X}) at its own position,
which preserves the reading frame of every downstream codon. Melting
temperature is estimated by the Wallace rule for short oligonucleotides and
by SantaLucia's nearest-neighbour thermodynamics with salt correction
\citep{santalucia1998} for longer duplexes within the range for which that
two-state model is applicable; the method and the primer/salt conditions
used are always reported alongside the value, and a duplex outside the
model's validated length range carries an explicit caveat rather than a bare
number.

\paragraph{Protein analysis.}
Protein descriptors include average molecular weight (IUPAC average residue
masses, less one water molecule per peptide bond), grand average of
hydropathy (GRAVY) and a Kyte--Doolittle hydropathy profile
\citep{kytedoolittle1982}, aromaticity and aliphatic index
\citep{ikai1980aliphatic}, extinction coefficient at 280\,nm by the
Gill--von Hippel method \citep{gillvonhippel1989} (reported both with
cysteines reduced and as cystines), and isoelectric point by iterative
charge-balance search over a Bjellqvist-style pKa table. Because independent
isoelectric-point calculators are documented to disagree by several tenths
of a pH unit depending on which of several published pKa scales they use
\citep{kozlowski2016ipc}, \bsi\ does not claim numerical identity with any
one external tool; instead its cross-validation suite
(Section~\ref{sec:cross-validation}) checks agreement with Biopython within
a stated, literature-consistent tolerance and, independently of any external
tool, checks that the net charge computed at the reported isoelectric point
is itself zero. A Chou--Fasman-style secondary-structure propensity sketch
\citep{choufasman1978} is provided for illustrative and teaching purposes
and is labelled, in the interface, the exported data, and this manuscript,
as a propensity sketch rather than a structure prediction; a user who needs
an actual secondary-structure assignment is directed to DSSP applied to a
real structure, or to a dedicated modern predictor.

\subsection{External structure-resource retrieval}
\label{sec:retrieval}

Three structural resources are queried, each through a dedicated provider
class implementing a common interface (Figure~\ref{fig:architecture}):

\begin{itemize}[leftmargin=1.4em]
  \item \textbf{RCSB PDB} \citep{berman2000pdb}: direct retrieval by PDB
  identifier, or sequence search followed by download of the top-scoring
  entry, with cross-reference accessions read from the matched polymer
  entity's reference-sequence metadata.
  \item \textbf{AlphaFold Protein Structure Database}
  \citep{varadi2022alphafolddb}: retrieval by UniProt accession. The
  provider queries the public prediction API to discover the model file
  currently served for that accession, rather than constructing the file
  URL from an assumed model-version number; the database periodically
  republishes its models under a new version, and a client that instead
  encodes a fixed version number degrades to reporting a coverage gap for
  every accession once the database moves past it, a failure mode that is
  indistinguishable from genuinely absent coverage without querying the
  authoritative endpoint. An HTTP 404 response from the prediction API is
  treated as a true coverage gap and is not retried; a malformed or empty
  API response is treated as a service failure and is.
  \item \textbf{ESM Atlas} (ESMFold, \citealp{lin2023esmfold}): on-demand
  structure prediction from an arbitrary sequence, subject to a configurable
  maximum sequence length; sequences exceeding the limit are rejected before
  any network request is made, and this rejection is reported as its own
  outcome category distinct from either a retrieved structure or a service
  failure (Section~\ref{sec:evaluation}).
\end{itemize}

UniProt \citep{uniprot2023} is queried first when the input is an
accession, to obtain the canonical sequence against which every retrieved
structure is subsequently validated; without this step, an accession query
would have no query sequence for the validation check of
Section~\ref{sec:validation} to use.

Providers are tried in a configurable order (RCSB, then AlphaFold DB, then
ESM Atlas, by default) and the manager stops at the first exact (M4,
Section~\ref{sec:m0m4}) match, unless configured to query every resource
regardless (used for the benchmark in Section~\ref{sec:live-benchmark}).
Responses are cached on disk, keyed by resource, query, and relevant
parameters, with a stored SHA-256 checksum verified before a cached payload
is served; a tampered or corrupted cache entry is discarded and refetched
rather than served.

\subsection{Sequence-to-structure validation}
\label{sec:validation}

Every structure a provider returns is compared against the query sequence
before being reported. The coordinate file is parsed to recover per-chain
sequence (from \texttt{SEQRES} records where present, falling back to the
$\text{C}\alpha$ trace otherwise), and the chain whose sequence aligns best
to the query is selected.

Alignment uses semi-global (``glocal'') dynamic programming rather than a
strict global alignment: the query is aligned end to end while the subject
chain may be a sub-sequence of it with no length penalty. Scoring is
deliberately simple (match $+1$, mismatch $-1$, linear gap $-2$, no
substitution matrix): the alignment's role here is to measure identity and
coverage of one candidate against the query, not to rank distant homologues
by evolutionary score. Because identity is taken over the aligned region only, it
depends on the scoring scheme: with permissive gap scoring, an alignment of two
unrelated sequences can reach a high identity over a short aligned stretch
(Section~\ref{sec:threshold}), so identity is only interpretable together with
coverage and, for the strongest claim, with the accession check. The algorithm has $O(mn)$ time and memory bound guarded by a maximum
cell-count cap that raises an explicit error rather than exhausting memory
on pathologically long inputs. This choice matters concretely: a deposited
construct that omits a cleaved signal peptide or a disordered terminal
region -- common in crystallography -- is still a correct structure of the
folded domain it does cover, and a strict global alignment
\citep{needleman1970} would understate its identity in proportion to the
omitted length. Two measures are reported: \emph{identity}, the fraction of
aligned positions that match, and \emph{coverage}, the fraction of the query
length included in the alignment; a fragment therefore shows high identity
with reduced coverage, which is the correct characterisation, rather than
being scored as a poor match overall.

Accessions are collected from the coordinate file's own cross-references
(\texttt{DBREF} records and equivalent \texttt{REMARK} fields) and, where
the originating provider's API exposes it, from provider-supplied metadata
independent of the coordinate file itself -- for RCSB, the matched polymer
entity's reference-sequence identifiers. The queried accession is
normalised (handling FASTA header forms such as
\texttt{sp\textbar P24941\textbar CDK2\_HUMAN} and AlphaFold model names such
as \texttt{AF-P69905-F1}) and compared against this set.

\subsection{M0--M4 sequence-to-structure correspondence classification}
\label{sec:m0m4}

The identity, coverage, and accession-match measures of
Section~\ref{sec:validation} are combined into a single ordinal
classification, summarised in Table~\ref{tab:m0m4} and
Figure~\ref{fig:decision}, reported alongside every retrieved structure in
both interfaces and in every exported table.

\begin{table}[htbp]
\centering
\small
\caption{The M0--M4 sequence-to-structure correspondence scale. Default
thresholds are shown; all six numeric thresholds are configurable and the
values in force are recorded with every result.}
\label{tab:m0m4}
\begin{tabular}{@{}llp{7.4cm}@{}}
\toprule
\textbf{Level} & \textbf{Name} & \textbf{Criteria} \\
\midrule
\textcolor{mapM4}{\mlevel{M4}} & Exact &
Accession cross-referenced by the structure \emph{and} identity $\geq 99\%$
\emph{and} coverage $\geq 95\%$. \\[2pt]
\textcolor{mapM3}{\mlevel{M3}} & High confidence &
Identity $\geq 95\%$ and coverage $\geq 80\%$; accession not confirmed
(the same or a near-identical sequence). \\[2pt]
\textcolor{mapM2}{\mlevel{M2}} & Related record &
Identity $\geq 40\%$ and coverage $\geq 40\%$, below the M3 thresholds
(a homologue, an orthologue, or a fragment); \emph{not} evidence of homology
(Section~\ref{sec:threshold}). \\[2pt]
\textcolor{mapM1}{\mlevel{M1}} & Unverified &
Either no comparison was possible (no query sequence, or no readable chain),
\emph{or} a comparison was made and fell below even the M2 thresholds. \\[2pt]
\textcolor{mapM0}{\mlevel{M0}} & No result &
No structural record was returned by any queried resource. \\
\bottomrule
\end{tabular}
\end{table}

\begin{figure}[htbp]
\centering
\begin{tikzpicture}[
    dec/.style={diamond, draw=black!70, aspect=2.4, minimum width=3.1cm, align=center, font=\scriptsize, inner sep=1pt},
    res/.style={rectangle, rounded corners, draw=black!70, minimum width=2.5cm, minimum height=0.75cm, align=center, font=\scriptsize\bfseries},
    lbl/.style={font=\scriptsize, midway, fill=white, inner sep=1pt}
  ]
  \node[res, fill=gray!15] (start) {Structure retrieved};
  \node[dec, below=0.55cm of start] (d1) {Comparison\\possible?};
  \node[res, right=1.65cm of d1, fill=mapM1!25, text=black] (m1a) {\textcolor{mapM1}{\textbf{M1}} unverified};
  \node[dec, below=0.85cm of d1] (d2) {identity $\geq$40\%\\coverage $\geq$40\%?};
  \node[res, right=1.65cm of d2, fill=mapM1!25, text=black] (m1b) {\textcolor{mapM1}{\textbf{M1}} unverified};
  \node[dec, below=0.85cm of d2] (d3) {identity $\geq$95\%\\coverage $\geq$80\%?};
  \node[res, right=1.65cm of d3, fill=mapM2!25, text=black] (m2) {\textcolor{mapM2}{\textbf{M2}} related record};
  \node[dec, below=0.85cm of d3] (d4) {accession\\cross-referenced\\\emph{and} identity $\geq$99\%,\\coverage $\geq$95\%?};
  \node[res, right=1.65cm of d4, fill=mapM3!25, text=black] (m3) {\textcolor{mapM3}{\textbf{M3}} high confidence};
  \node[res, below=0.85cm of d4, fill=mapM4!25, text=black] (m4) {\textcolor{mapM4}{\textbf{M4}} exact};

  \draw[-{Stealth[length=1.8mm]}] (start) -- (d1);
  \draw[-{Stealth[length=1.8mm]}] (d1) -- node[lbl]{no} (m1a);
  \draw[-{Stealth[length=1.8mm]}] (d1) -- node[lbl, anchor=west]{yes} (d2);
  \draw[-{Stealth[length=1.8mm]}] (d2) -- node[lbl]{no} (m1b);
  \draw[-{Stealth[length=1.8mm]}] (d2) -- node[lbl, anchor=west]{yes} (d3);
  \draw[-{Stealth[length=1.8mm]}] (d3) -- node[lbl]{no} (m2);
  \draw[-{Stealth[length=1.8mm]}] (d3) -- node[lbl, anchor=west]{yes} (d4);
  \draw[-{Stealth[length=1.8mm]}] (d4) -- node[lbl]{no} (m3);
  \draw[-{Stealth[length=1.8mm]}] (d4) -- node[lbl, anchor=west]{yes} (m4);
\end{tikzpicture}
\caption{Decision path from a retrieved structure to its M0--M4
classification. M0 (no structural record returned by any resource) is not
shown, as it precedes this diagram entirely. Only M4 supports an unqualified
statement that a retrieved structure is the queried protein.}
\label{fig:decision}
\end{figure}

Two consequences of this design are worth stating explicitly. First, accession agreement alone does not confer M4: a
short peptide fragment of a much longer protein can legitimately carry its
parent accession while covering a small fraction of the query, so M4
additionally requires the coverage threshold; without it, a fragment would
be indistinguishable from a complete structure. Second, sequence agreement
alone does not confer M4 either: two distinct proteins can share a
near-identical local region, and a deposited construct may be a chimera or
engineered variant whose \emph{chain sequence} matches the query while the
\emph{entry} describes something else; a perfect sequence match without
accession confirmation is therefore reported as M3, not M4, and the
difference between ``very likely correct'' and ``confirmed correct'' is
preserved in the report rather than collapsed. A structure whose measured
identity and coverage fall below even the M2 thresholds is classified M1
with an explicit warning naming the accession actually returned, rather than
being rounded up to M2 (``related'').

\paragraph{What ``accession confirmed'' means, and does not mean.} The
accession check reads the structure's own cross-references (coordinate-file
\texttt{DBREF} records and, for RCSB, the matched entity's reference-sequence
identifiers). It therefore asserts that the \emph{entry} is catalogued under the
queried accession, which is an accession-level claim: an identical sequence
filed under a different accession is not confirmed (it is reported M3), and the
check inherits whatever curation errors the cross-references contain. The
validation study (Section~\ref{sec:validation-study}) judges correctness at this
same accession level, against references that are curated independently of this
software but are not independent of the cross-references the check reads.

\paragraph{Partial-structure qualifier.} A deposited structure is often only a
domain construct. When the accession is confirmed and identity is at least the
high-identity threshold ($95\%$) but coverage is below the high-coverage
threshold ($80\%$), the report carries \texttt{partial\_structure = true} and a
message stating that the structure is a real but incomplete structure of the
requested protein. The qualifier does not change the mapping level (such a
result is M1 or M2) and cannot co-occur with M3 or M4, so M4 keeps its meaning
of a full-length exact match. It was added because the validation study showed
that the original message for these results (``does not correspond to the
requested protein'') was wrong for them.

\subsection{Error handling and fault tolerance}
\label{sec:fault-tolerance-methods}

HTTP requests are wrapped in a typed error hierarchy rooted at a common
\texttt{BioSeqInsightError} base class, with specific subclasses
distinguishing a transport failure, a timeout, a not-found result, a
service error, a validation failure, and a cache failure, so that a caller
can handle each category differently. Transient
failures (HTTP 429, 500, 502, 503, 504; connection errors; read timeouts)
are retried with exponential backoff and random jitter up to a configurable
attempt budget, honouring a server-supplied \texttt{Retry-After} header when
present; permanent client errors (404, 403) are not retried. Every attempt
is recorded (status code or transport error, elapsed time, whether a retry
followed) and the full attempt history is exposed in every result, which is
what makes the fault-injection evaluation of
Section~\ref{sec:fault-injection} directly traceable to specific,
inspectable request sequences rather than to an aggregate pass/fail count
alone.

A response is additionally validated as \emph{content} before being
accepted, independent of its HTTP status: a payload that does not contain an
expected marker (for example, an \texttt{ATOM} record in a coordinate-file
response) is treated as a failure and retried according to the same policy,
even when the server returned HTTP 200. This addresses a specific,
concrete failure pattern in which a gateway or load balancer returns an
HTML error page with a 200 status during an outage; a client that trusts
the status code alone would write that page to disk as though it were a
structure file.

When the first-tried resource cannot supply a usable result, the manager
proceeds to the next configured resource rather than reporting failure
immediately (Section~\ref{sec:retrieval}); every resource that was queried,
successfully or not, is reported in the result, so a fallback is visible
rather than silent.

\subsection{Reproducibility and project workflows}
\label{sec:reproducibility}

Configuration (retry budget and backoff parameters, cache location and
time-to-live, provider fallback order, M0--M4 thresholds, melting-temperature
conditions, default genetic-code table, batch worker count, and endpoint
URLs) resolves through a documented precedence chain -- built-in defaults,
then an optional settings file, then environment variables, then explicit
command-line or API arguments -- and is validated on load, with inconsistent
values (for example, an M2 identity threshold exceeding the M3 threshold, or
an unrecognised provider name) rejected immediately rather than producing an
unexplained result later. The complete effective settings, together with the
software version, are recorded with every batch run and every project.

A \emph{project} is a directory holding the exact input sequences analysed,
any retrieved structure coordinates, exported result tables, structured logs,
and a metadata file recording the software version and settings snapshot
above. A project exports to a single zip archive that re-imports into a
working project on another machine, with archive extraction defended against
path-traversal (``zip slip'') by rejecting any entry whose path would
escape the destination directory. This is the mechanism by which a
completed analysis in \bsi\ can be handed to a collaborator, or re-run
against a later software version, as a single, self-contained, auditable
unit rather than as an ad hoc collection of files and remembered settings;
its effect on the number of manual steps this requires is quantified in
Section~\ref{sec:workflow-comparison}.
